\documentclass[12pt, letterpaper]{article}

\usepackage[margin=1in]{geometry}
\usepackage{booktabs}
\usepackage{array}
\usepackage{longtable}
\usepackage{amsmath}
\usepackage{hyperref}
\usepackage{xcolor}
\usepackage{parskip}
\usepackage{titlesec}
\usepackage{enumitem}
\usepackage{caption}
\usepackage{float}

\hypersetup{
    colorlinks=true,
    linkcolor=black,
    urlcolor=blue,
    citecolor=black
}

\titleformat{\section}{\large\bfseries}{}{0em}{}[\titlerule]
\titleformat{\subsection}{\normalsize\bfseries}{}{0em}{}

\title{\textbf{Algorithmic Fairness Audit} \\[0.4em]
       \large Student Exam Score Prediction \\
       Kaggle Playground Series S6E1}
\author{George Buck \and Sinclair Madden}
\date{May 2026}

\begin{document}

\maketitle
\tableofcontents
\newpage

% ─────────────────────────────────────────────────────────────────────────────
\section{Dataset}

\subsection{Source and Size}

The dataset is drawn from Kaggle Playground Series Season 6, Episode 1 (S6E1),
a synthetic tabular regression competition centred on student academic
performance. The training set contains \textbf{630,000 rows} and \textbf{11
features}; the target variable is \texttt{exam\_score}, a continuous value on
the interval $[0, 100]$.

\subsection{Features}

\begin{table}[H]
\centering
\caption{Feature descriptions}
\begin{tabular}{@{}llp{7cm}@{}}
\toprule
\textbf{Feature} & \textbf{Type} & \textbf{Description} \\
\midrule
\texttt{age}               & Numeric      & Student age (17--24) \\
\texttt{study\_hours}      & Numeric      & Average daily study hours \\
\texttt{class\_attendance} & Numeric      & Percentage of classes attended \\
\texttt{sleep\_hours}      & Numeric      & Average nightly sleep hours \\
\texttt{gender}            & Categorical  & Self-identified gender (female, male, other) \\
\texttt{course}            & Categorical  & Enrolled academic programme \\
\texttt{internet\_access}  & Categorical  & Reliable internet connection (yes/no) \\
\texttt{sleep\_quality}    & Categorical  & Self-reported sleep quality \\
\texttt{study\_method}     & Categorical  & Primary study method \\
\texttt{facility\_rating}  & Categorical  & Student rating of campus facilities \\
\texttt{exam\_difficulty}  & Categorical  & Self-reported exam difficulty \\
\bottomrule
\end{tabular}
\end{table}

\subsection{Descriptive Statistics}

\begin{table}[H]
\centering
\caption{Key numeric feature statistics (training set, $n = 630{,}000$)}
\begin{tabular}{@{}lrrrr@{}}
\toprule
\textbf{Feature} & \textbf{Mean} & \textbf{Std} & \textbf{Min} & \textbf{Max} \\
\midrule
\texttt{age}               & 20.55 & 2.26 & 17.0 & 24.0 \\
\texttt{study\_hours}      &  4.00 & 2.36 &  0.1 & 14.0 \\
\texttt{class\_attendance} & 72.0  & 17.4 & 40.6 & 100.0 \\
\texttt{sleep\_hours}      &  7.07 & 1.74 &  4.1 & 12.0 \\
\texttt{exam\_score}       & 62.5  & 18.9 &  0.0 & 100.0 \\
\bottomrule
\end{tabular}
\end{table}

Notable distributional facts:
\begin{itemize}[noitemsep]
    \item \textbf{Internet access:} 92\% of students report having access; 8\% do not.
    \item \textbf{Gender:} Roughly balanced across female ($\approx$33\%), male ($\approx$33\%), and other ($\approx$33\%) in both training and test splits.
    \item \textbf{Study hours} is the single numeric feature most correlated with exam score ($r = 0.762$); class attendance is second ($r = 0.361$). Age is essentially uncorrelated ($r = 0.010$).
\end{itemize}

% ─────────────────────────────────────────────────────────────────────────────
\section{Methods}

\subsection{Original Model (Audited)}

The model under audit is a \textbf{Histogram-Based Gradient Boosting Regressor}
(\texttt{HistGradientBoostingRegressor}) from \texttt{scikit-learn}, trained
with the following hyperparameters:

\begin{center}
\texttt{max\_iter=1000,\ learning\_rate=0.05,\ max\_depth=10,\ random\_state=42}
\end{center}

Pre-processing used \texttt{LabelEncoder} (alphabetical mapping) applied
independently to each of the seven categorical columns. The data were split
80/20 into training and test sets (\texttt{random\_state=42}). \textbf{No
changes were made to the model or its training pipeline during this audit.}

\subsection{Exploratory Data Analysis and Statistical Tests}

Prior to the fairness audit, the following statistical tests were applied to the
raw training data:

\begin{itemize}[noitemsep]
    \item \textbf{Pairwise correlation heatmap} across all numeric features.
    \item \textbf{Independent-samples $t$-test:} internet access group vs.\ exam score.
    \item \textbf{One-way ANOVA:} study method vs.\ exam score, with \textbf{Tukey HSD} post-hoc comparisons across all method pairs.
    \item \textbf{Chi-square test of independence:} gender $\times$ exam difficulty.
    \item \textbf{OLS regression:} numeric features regressed on exam score as a linear baseline.
\end{itemize}

\subsection{Fairness Audit}

Seven fairness tests were applied to the frozen, pre-trained model using the held-out
test split ($n = 126{,}000$). The pass/fail threshold for binary classification tasks
was set at the training-set mean exam score (62.5). Protected attributes evaluated
were \texttt{gender} and \texttt{internet\_access} (as a socioeconomic proxy).

\begin{table}[H]
\centering
\caption{Fairness tests applied}
\begin{tabular}{@{}clp{8cm}@{}}
\toprule
\textbf{\#} & \textbf{Test} & \textbf{What it measures} \\
\midrule
1 & Data-level bias         & Are raw score distributions skewed across gender groups? \\
2 & Error parity            & Are MAE, RMSE, and prediction bias equal across groups? \\
3 & Demographic parity      & Do all groups receive the same predicted pass rate? \\
4 & Equal opportunity / equalized odds & Are TPR and FPR equal across groups? \\
5 & Counterfactual fairness & Does flipping gender alone change predictions? \\
6 & Permutation feature importance & Does the model rely on gender as a predictor? \\
7 & Socioeconomic proxy     & Does internet access predict higher predicted scores? \\
\bottomrule
\end{tabular}
\end{table}

\subsection{SHAP Analysis}

\texttt{shap.TreeExplainer} was used to compute exact Shapley values for a
5,000-row random sample of the test set (\texttt{random\_state=7}). Because
\texttt{HistGradientBoostingRegressor} is a gradient-boosted tree ensemble,
\texttt{TreeExplainer} produces exact (non-approximated) decompositions.
Results were visualised as global bar charts, beeswarm plots, and dependence
plots for the two strongest features. A helper function,
\texttt{explain\_student()}, provides per-student prediction explanations.

% ─────────────────────────────────────────────────────────────────────────────
\section{Outcomes}

\subsection{Model Performance}

\begin{center}
\begin{tabular}{@{}lr@{}}
\toprule
\textbf{Metric} & \textbf{Value} \\
\midrule
$R^2$ (test set) & 0.7836 \\
MAE (test set)   & 6.998 \\
\bottomrule
\end{tabular}
\end{center}

\subsection{EDA and Statistical Test Results}

\begin{itemize}[noitemsep]
    \item \textbf{Internet access vs.\ exam score:} The independent-samples $t$-test
          returned $t = 0.355$, $p = 0.723$, indicating no statistically significant
          difference in mean exam scores between students with and without internet access.
    \item \textbf{Study method vs.\ exam score:} One-way ANOVA returned $F = 8{,}304$,
          $p \approx 0$. Tukey HSD confirmed that every pair of study methods differed
          significantly ($\alpha = 0.05$), with self-study associated with the lowest
          mean scores and coaching with the highest.
    \item \textbf{Gender $\times$ exam difficulty:} The chi-square test returned
          $\chi^2 = 40.8$, $p = 2.9 \times 10^{-8}$. Although statistically significant
          at this sample size, the effect size is negligible given $n = 630{,}000$.
\end{itemize}

\subsection{Fairness Audit Results}

\begin{table}[H]
\centering
\caption{Fairness audit summary}
\begin{tabular}{@{}lllc@{}}
\toprule
\textbf{Metric} & \textbf{Value} & \textbf{Benchmark} & \textbf{Pass} \\
\midrule
Gender MAE spread (max $-$ min)              & 0.023   & $\approx 0$ is fair        & \checkmark \\
Gender bias spread (max $-$ min)             & 0.062   & $\approx 0$ is fair        & \checkmark \\
Demographic parity diff (male $-$ female)    & $-$0.003 & $|$diff$|$ $< 0.05$       & \checkmark \\
Equal opportunity diff (male $-$ female)     & $-$0.005 & $|$diff$|$ $< 0.05$       & \checkmark \\
Counterfactual mean $|\Delta|$ (gender flip) & 0.067   & 0 = perfectly fair         & \checkmark \\
Predictions shifted $>1$ pt by gender flip   & 0.0\%   & ---                        & \checkmark \\
Gender permutation importance rank           & 9th/11  & Low rank = fair            & \checkmark \\
\bottomrule
\end{tabular}
\end{table}

\paragraph{Error parity by gender.}
MAE varied by at most 0.023 points across gender groups (female: 7.00, male: 6.99,
other: 7.01). Prediction bias was similarly uniform (female: $-$0.04, male: $-$0.05,
other: $+$0.01). Both spreads are negligible relative to the model's overall MAE of 6.998.

\paragraph{Demographic parity.}
Predicted pass rates were 50.7\% (female), 50.4\% (male), and 51.0\% (other), against
actual rates of 50.2\%, 50.4\%, and 50.6\% respectively. The maximum demographic parity
difference was 0.006, well within the 0.05 benchmark of Feldman et al.\ (2015).

\paragraph{Equal opportunity and equalized odds.}
TPR ranged from 0.863 (male) to 0.874 (other); FPR ranged from 0.137 (other) to
0.142 (female). The maximum equal opportunity difference across any pair of groups
was 0.011, well within the 0.05 benchmark of Hardt, Price, and Srebro (2016).

\paragraph{Counterfactual fairness.}
Flipping gender from female to male (holding all other features fixed, $n = 2{,}000$)
produced a mean prediction shift of $+$0.001, a mean absolute shift of 0.067, and a
maximum absolute shift of 0.900. No prediction changed by more than 1 point.

\paragraph{Socioeconomic proxy.}
The actual exam score gap between students with and without internet access was
$+$0.062 points (yes: 62.61, no: 62.55). The predicted gap was $-$0.059 points.
The reported gap amplification factor of $-$0.961 is numerically unreliable: with
an actual gap of only $\approx 0.06$ points, the ratio \texttt{pred\_gap / actual\_gap}
is highly sensitive to near-zero denominators. The substantive finding is that both
actual and predicted gaps are negligibly small in absolute terms, consistent with
the $t$-test result ($p = 0.723$).

\subsection{SHAP Feature Importance}

\begin{table}[H]
\centering
\caption{Permutation importance and mean $|$SHAP$|$ (ranked)}
\begin{tabular}{@{}lrl@{}}
\toprule
\textbf{Feature} & \textbf{Perm.\ importance ($\Delta R^2$)} & \textbf{Note} \\
\midrule
\texttt{study\_hours}      & 1.030  & Primary driver \\
\texttt{class\_attendance} & 0.165  & \\
\texttt{sleep\_quality}    & 0.077  & \\
\texttt{study\_method}     & 0.067  & \\
\texttt{facility\_rating}  & 0.047  & \\
\texttt{sleep\_hours}      & 0.030  & \\
\texttt{age}               & 0.001  & \\
\texttt{course}            & 0.0003 & \\
\texttt{gender}            & 0.0002 & Protected attribute \\
\texttt{exam\_difficulty}  & 0.0000 & \\
\texttt{internet\_access}  & 0.0000 & Socioeconomic proxy \\
\bottomrule
\end{tabular}
\end{table}

SHAP beeswarm and dependence plots confirm that \texttt{study\_hours} dominates
model predictions: higher study hours are associated with large positive SHAP values,
with the effect non-linear at the extremes. \texttt{class\_attendance} is the second
most influential feature. Both \texttt{gender} and \texttt{internet\_access} have
mean $|$SHAP$|$ values indistinguishable from zero, confirming they contribute
negligibly to any individual prediction.

\subsection{Overall Assessment}

The model shows \textbf{no meaningful disparity} across gender groups on any
fairness criterion tested. Both protected attributes (\texttt{gender} and
\texttt{internet\_access}) rank last in permutation importance and SHAP magnitude.
Counterfactual gender flips produce negligible prediction changes. Demographic
parity and equal opportunity differences are an order of magnitude below standard
fairness thresholds. The dominant driver of predicted exam score is
\texttt{study\_hours}, followed by \texttt{class\_attendance} --- both
effort-based, non-demographic features.

% ─────────────────────────────────────────────────────────────────────────────
\section*{References}

\begin{itemize}[noitemsep]
    \item Feldman, M., Friedler, S., Moeller, J., Scheidegger, C., \& Venkatasubramanian, S. (2015). Certifying and removing disparate impact. \textit{Proceedings of KDD 2015}.
    \item Hardt, M., Price, E., \& Srebro, N. (2016). Equality of opportunity in supervised learning. \textit{Advances in Neural Information Processing Systems (NeurIPS) 2016}.
\end{itemize}

\end{document}
